\documentclass[10pt,aspectratio=169,serif]{beamer}

\usepackage[fontset=none, scheme=plain]{ctex}
\IfFontExistsTF{Songti SC}{%
	\setCJKmainfont{Songti SC}[AutoFakeSlant=0.2]%
	\setCJKsansfont{Songti SC}[AutoFakeSlant=0.2]%
}{%
	\setCJKmainfont{FandolSong-Regular}[Extension=.otf,BoldFont=FandolSong-Bold,AutoFakeSlant=0.2]%
	\setCJKsansfont{FandolSong-Regular}[Extension=.otf,BoldFont=FandolSong-Bold]%
}
\usepackage{amsmath,amssymb,graphicx,tikz,booktabs,multicol,calligra}
\usepackage{pgfplots}
\pgfplotsset{compat=1.18}
\usepackage{tcolorbox}
\tcbuselibrary{skins}
\usepackage{XyygU}
\usepackage{caption}
\captionsetup[figure]{name=图,labelsep=period}

\usetikzlibrary{arrows.meta,positioning,calc}

\setlength{\abovedisplayskip}{0pt plus 2pt minus 2pt}
\setlength{\belowdisplayskip}{0pt plus 2pt minus 2pt}

\title[函数的概念练习]{3.1.1 函数的概念（练习）}
\subtitle{定义域 \& 函数值 \& 同一函数}
\author{}
\institute{}
\date{2026-10-16}

\begin{document}

\begin{frame}
  \titlepage
\end{frame}

% ========== 第1题 ==========
\begin{frame}[t]{第1题：求函数的定义域}
\textbf{1.} 求下列函数的定义域：

\vspace{1em}
(1) $\displaystyle f(x)=\frac{1}{4x+7}$；

\vspace{1.5em}
(2) $\displaystyle f(x)=\sqrt{1-x}+\sqrt{x+3}-1$。
\end{frame}

\begin{frame}[t]{第1题 解答}
\begin{tcolorbox}[colback=orange!10,colframe=orange,title=\textbf{解}]
(1) 要使分式 $\dfrac{1}{4x+7}$ 有意义，分母不能为零：
\[
4x+7 \neq 0 \quad\Longrightarrow\quad x \neq -\frac{7}{4}.
\]
\textbf{定义域：} $\left(-\infty,\,-\dfrac{7}{4}\right)\cup\left(-\dfrac{7}{4},\,+\infty\right)$。

\vspace{1em}
(2) 要使两个根式都有意义，需同时满足：
\[
\begin{cases}
1-x \ge 0,\\
x+3 \ge 0,
\end{cases}
\quad\Longrightarrow\quad
\begin{cases}
x \le 1,\\
x \ge -3.
\end{cases}
\]
取交集得 $-3 \le x \le 1$。

\textbf{定义域：} $[-3,\,1]$。
\end{tcolorbox}
\end{frame}

% ========== 第2题 ==========
\begin{frame}[t]{第2题：求函数值}
\textbf{2.} 已知函数 $f(x)=3x^3+2x$，

\vspace{1em}
(1) 求 $f(2)$，$f(-2)$，$f(2)+f(-2)$ 的值；

\vspace{1em}
(2) 求 $f(a)$，$f(-a)$，$f(a)+f(-a)$ 的值。
\end{frame}

\begin{frame}[t]{第2题 解答}
\begin{tcolorbox}[colback=orange!10,colframe=orange,title=\textbf{解}]
(1) 分别代入：
\[
f(2)=3\cdot 2^3+2\cdot 2=24+4=28,
\]
\[
f(-2)=3\cdot(-2)^3+2\cdot(-2)=-24-4=-28,
\]
\[
f(2)+f(-2)=28+(-28)=0.
\]

\vspace{0.5em}
(2) 一般情形：
\[
f(a)=3a^3+2a,
\]
\[
f(-a)=3(-a)^3+2(-a)=-3a^3-2a=-(3a^3+2a),
\]
\[
f(a)+f(-a)=(3a^3+2a)+(-3a^3-2a)=0.
\]

\vspace{0.3em}
\small（这说明 $f(x)=3x^3+2x$ 是一个奇函数。）
\end{tcolorbox}
\end{frame}

% ========== 第3题 ==========
\begin{frame}[t]{第3题：判断是否为同一函数}
\textbf{3.} 判断下列各组中的函数是否为同一个函数，并说明理由：

\vspace{1em}
(1) 表示炮弹飞行高度 $h$ 与时间 $t$ 关系的函数 $h=130t-5t^2$ 和二次函数 $y=130x-5x^2$；

\vspace{1em}
(2) $f(x)=1$ 和 $g(x)=x^0$。
\end{frame}

\begin{frame}[t]{第3题 解答}
\begin{tcolorbox}[colback=orange!10,colframe=orange,title=\textbf{解}]
两个函数是同一个函数，当且仅当 \textbf{定义域相同} 且 \textbf{对应关系完全一致}。

\vspace{0.5em}
(1) \alert{不是}同一个函数。

炮弹飞行高度 $h=130t-5t^2$ 中，时间 $t$ 受实际意义限制：从发射到落地，$0 \le t \le 26$（令 $h=0$ 得 $t=0$ 或 $26$），故定义域为 $[0,26]$；而二次函数 $y=130x-5x^2$ 的定义域为 $\mathbf{R}$。两者定义域不同，故不是同一函数。

\vspace{0.5em}
(2) \alert{不是}同一个函数。

$f(x)=1$ 的定义域为 $\mathbf{R}$；而 $g(x)=x^0=1$ 要求 $x\neq 0$，定义域为 $(-\infty,0)\cup(0,+\infty)$。两者定义域不同，故不是同一函数。
\end{tcolorbox}
\end{frame}

\begin{frame}[c]
  \begin{center}
    {\Huge\calligra Thanks!}
  \end{center}
\end{frame}

\end{document}
